\begin{otherlanguage*}{english}
\nocite{Grothendieck1990XVIII}
The opening of Chapter XVIII asks what a category of coefficients must retain if the formal operations of sheaf theory are to survive outside a topos. Grothendieck begins with colimits and a small dense collection of test objects. The proposed tensor then forces him to ask whether continuity conditions define a reflective category, and that question leads to accessibility. The long investigation of cardinal filtrations is therefore part of the construction of coefficients. His letter to Thomason of 2 April 1991 explicitly describes this chapter as a return to the accessibility ideas of SGA~4, Expos\'e I, Section 9 \citep{Grothendieck1991Letter,SGA4I1972}.

These pages run from manuscript page 50 to page 138, corresponding to printed pages 1--43 of the edition. They contain Sections 4.1--4.5: pseudo-topoi and the tensor problem, its small-presheaf formulation, the proposed accessibility criterion, the restricted ind-completion, and the intrinsic cardinal filtration. The text often records an unsuccessful proof and then retains the objection which caused its abandonment. A faithful reading has to preserve this exploratory structure. An assertion subsequently withdrawn is not evidence that the intended theorem is false; conversely, the truth of the final theorem does not validate an intermediate equivalence.

We work in a fixed Grothendieck universe and use smallness relative to that universe. One translation governs the entire discussion. If $\pi$ is infinite, a category called \emph{grande devant $\pi$} is filtered with respect to diagrams having at most $\pi$ arrows. In the modern regular-cardinal convention it is $\pi^+$-filtered. Grothendieck's $\pi$-accessible objects are consequently $\pi^+$-presentable objects. His cocomplete categories with accessible small generators are locally presentable categories; the broader modern notion of an accessible category does not include cocompleteness. The cardinal attached to an object and one uniform cardinal for a chosen small generating family must also be distinguished.

A pseudo-topos has enough density for the dual special adjoint functor theorem. Colimit-preserving functors out of it have right adjoints, and a contravariant set-valued functor taking colimits to limits is representable. These facts make the opening construction plausible. They also expose a distinction in the motivating terminology. If morphisms of pseudo-topoi are described by their inverse-image functors, then the actual bicategorical product has underlying category $\mathcal M\times\mathcal N$: a cocontinuous functor out of that category is determined by its restrictions to $(m,0)$ and $(0,n)$. The proposed $\mathcal M\boxtimes\mathcal N$ represents functors cocontinuous in each variable separately. It is a tensor, with $\mathbf{Set}$ as unit, rather than this Cartesian product. For Grothendieck topoi the tensor nevertheless agrees with the product in the more restrictive category of topoi and geometric morphisms; Pitts's Theorem 2.3 proves the ordinary comparison \citep{Pitts1985}.

The tensor is first described as the category of bifunctors
\[
 \mathcal M^{\mathrm{op}}\times\mathcal N^{\mathrm{op}}
 \longrightarrow\mathbf{Set}
\]
which take colimits to limits separately. Density embeds it in a presheaf category on a small product of test categories. This proves local smallness and gives a useful formula for its continuous extensions, but it does not by itself provide the reflector. Nor does the naive external product of two representable presheaves usually satisfy the continuity condition: the empty colimit already forces both representing objects to be terminal. The example $\mathbf{Ab}\boxtimes\mathbf{Ab}\simeq\mathbf{Ab}$ makes the failure concrete. Its external tensor is the ordinary tensor product over $\mathbb Z$, and it cannot be recovered from the Cartesian product formula suggested in the manuscript.

Accessibility supplies the missing reflector. Choose a common regular cardinal and small full subcategories of presentable objects closed under colimits below that cardinal. Separate continuity can now be imposed by orthogonality to a \emph{set} of maps between small presheaves. The commentary constructs the reflection by alternately adjoining extensions and identifying extensions which agree on their domains. Presentability controls the length of the construction. The resulting category is again locally presentable, and reflection of representables supplies its elementary tensors. This gives a complete existence proof in the setting that the chapter is trying to isolate. Kelly's treatment of categories with prescribed colimits supplies the ordinary categorical background \citep[Section 6.5]{Kelly1982}. Lurie's tensor on presentable $\infty$-categories is a later comparison with the same universal pattern, not a substitute for the ordinary cardinal estimates here \citep[Section 4.8.1]{Lurie2017HA}.

The structural assertion behind Theorem 4.3.6 survives the failure of its first numerical proof. Starting from a small dense family of presentable objects, close that family under suitably small colimits. Every slice of the resulting category over an object is filtered at the required rank, and its canonical colimit is that object. The ind-object Hom formula proves full faithfulness. Thus the category is a restricted ind-completion. This argument does not require the provisional filtration by sizes of categories of elements to have the properties claimed for it on pages 72--76. The latter proof uses an unestablished directed-hull bound and imports a covering estimate from SGA~4 without its universality hypothesis. The manuscript's own corrective footnotes are justified.

Section 4.4 gives the complementary construction. When $C$ has colimits of diagrams with at most $\pi$ arrows, its $\pi^+$-ind-completion is cocomplete and is the free cocomplete category which preserves those existing colimits. The proof passes through cofinality, simultaneous presentations of diagrams, bounded coproducts, and the expression of an arbitrary indexing category as a directed union of small subcategories. We supply the omitted comma-adjunction argument and separate the existence of individual ind-presentations from closure under filtered colimits. The unrestricted claim that restriction on all functor categories is fully faithful is false; the needed restriction to functors preserving the density presentations is fully faithful and repairs the universal-property proof. Henry's recent theorem gives an exact modern comparison for these diagram presentations, with the hypotheses on the diagram shape and on $C$ stated explicitly \citep{Henry2025}.

The difference between ordinary and more highly filtered ind-objects is essential. Let $C$ consist of countable sets. Its $\aleph_1$-ind-completion is sets, but the ordinary ind-object obtained from finite subsets of $\mathbb Z$ has, at $c$, only the finite-image maps $c\to\mathbb Z$. Translation gives it a $\mathbb Z$-action which cannot occur at any nonempty equivariant representable stage. This answers negatively the specific cofinality question on manuscript page 94. It also explains why the restricted completion need not be closed under ordinary filtered colimits calculated in the presheaf category.

Section 4.5 measures an object in two ways. Write $\mathcal M_\pi$ for its intrinsic presentability layer, and $\Filt^\pi$ for the layer whose restricted nerve has at most $\pi$ maps from each test object, once the chosen small category itself is bounded by $\pi$. In sets, with tests of size at most $\pi_0$, these conditions read respectively
\[
 |X|\leq\pi,
 \qquad |X|^{\pi_0}\leq\pi.
\]
They need not coincide. The corrected comparison is
\[
 \Filt^\pi\subseteq\mathcal M_\pi
        \subseteq\Filt^{\pi^{\pi_0}}.
\]
The fixed-point condition $\pi^{\pi_0}=\pi$ restores equality. No continuum hypothesis is needed for the comparison itself. The counterexample sketched in the manuscript by a supposedly smallest directed family of subsets is not valid; the commentary replaces it with a diagonal argument at a singular cardinal of countable cofinality.

There is a second distinction within the proposed primed filtration. An object may be the colimit of a diagram of at most $\pi$ objects and arrows from $C$ without admitting a cofinal diagram of that size in $C/X$. For $C$ the countable sets and $|X|=\beth_\omega$, the coproduct of the singleton subsets gives the first presentation. A cofinal presentation would give at most $\beth_\omega$ countable subsets cofinal among all countable subsets of $X$, which the same diagonal argument excludes. Thus the two triples of conditions in Proposition 4.5.13 cannot be identified.

The positive outcome is stronger than the closing warning on page 138 suggests. The class of \emph{one-step} colimits of $C$-diagrams with at most $\pi$ arrows is already $\mathcal M_\pi$, for every $\pi\geq\pi_0$. The commentary proves closure under coequalizers by recording, for each generator in a presentation of the source, only the bounded portion of a target presentation through which its two images factor. This keeps the enlarged diagram within the original cardinal bound. Coproducts and splitting of idempotents finish the proof. The general presentation theorem is known; Barton's authored proof discusses its earlier attribution and the difficulty with the originally cited justification \citep{Barton2019Presentable}. Our argument is included as a direct repair in Grothendieck's setting, without a claim of priority. It also yields the sharper bounds on the number of objects and arrows which the chapter seeks.

Two later corrections prevent the descent arguments from obscuring that conclusion. Presentability factors an identity through a stage as a \emph{retract}; Lemma 4.5.14 must retain closure under retracts before replacing this by membership in the chosen subcategory. And a universal epimorphism need not transmit a bound on the number of relations. A category of sets with one countably-ary relation gives a counterexample to Corollary 4.5.16 even when all colimits are universal. Requiring the outer family to be universally \emph{strict} epimorphic restores the corollary. Both corrections leave the intrinsic filtration theorem intact.

Footnote 50 suggests testing universality of filtered colimits in the generating category. We prove a precise version: if $C$ also has pullbacks, then ordinary filtered colimits in $\Ind'_{\pi_0}(C)$ are universal exactly when the filtered colimits in $C$ with at most $\pi_0$ arrows are universal. The proof approximates the entire cocone and its base-change arrow together, and then expresses any filtered indexing category as a $\pi_0^+$-directed union of filtered subcategories with at most $\pi_0$ arrows. The additional pullback hypothesis makes the test intrinsic to $C$.

The remaining verification limits concern the disputed implication in Proposition 4.3.1 for arbitrary pseudo-topoi and the suggestion that universality of strict epimorphic families implies universality of ordinary filtered colimits. We do not describe these as literature-established open problems. The relevant older discussions of density, ind-objects, and the distinction between formal and realized colimits are cross-referenced where they are used.
\end{otherlanguage*}
